\begin{textsample}{NEG Dim 3 – ma – Score -2.00 – ma/ma\_em\_28.txt}  \label{ex:f3_neg_126}
Os objetivos dos itinerários formativos Para a organização dos itinerários formativos sugeridos para o território maranhense, considera -se os objetivos previstos nos Referenciais para Elaboração dos Itinerários Formativos, publicados na Portaria nº 1432, de 28 de dezembro de 2018, quais sejam : - Aprofundar as aprendizagens relacionadas às \textbf{competências} gerais, às áreas de conhecimento e/ou à formação técnica e profissional ; - Consolidar a formação integral dos estudantes, desenvolvendo a autonomia necessária para que realizem seus projetos de vida ; - Promover a incorporação de valores universais, como ética, liberdade, democracia, justiça social, pluralidade, solidariedade e sustentabilidade ; e - Desenvolver habilidades que permitam aos estudantes ter uma visão de mundo ampla e heterogênea, tomar decisões e agir nas mais diversas situações, seja na escola, seja no trabalho, seja na vida.
Para cumprir tais objetivos, os itinerários formativos se constituem a partir dos componentes curriculares da Base Nacional Comum Curricular e da parte de formação diversificada, já desenvolvidos nos Centros Educa Mais desde a sua implantação, tanto nas metodologias e práticas conhecidas, que sofreram alterações, como nos novos componentes concebidos para o Novo Ensino Médio pelo Instituto de Corresponsabilidade pela Educação.
Esses componentes curriculares se organizam a partir dos quatro eixos estruturantes⁶, que são complementares e coexistem nos distintos itinerários, garantindo que os estudantes experimentem situações de aprendizagem e vivências e possam desenvolver um conjunto diferenciado e ampliado de habilidades importantes e significativas para sua formação integral

% matched lemmas: competência
\end{textsample}
